\documentclass[11pt,letterpaper]{article}

\usepackage[margin=1in]{geometry}
\usepackage{amsmath,amssymb,amsthm}
\usepackage{physics}
\usepackage{booktabs}
\usepackage{natbib}
\usepackage[hyperfootnotes=false]{hyperref}
\usepackage{cleveref}
\usepackage{float}
\usepackage{titlesec}
\usepackage{microtype}

% Section formatting
\titleformat{\section}{\large\bfseries\uppercase}{\thesection}{1em}{}
\titleformat{\subsection}{\normalsize\bfseries}{\thesubsection}{1em}{}

\hypersetup{
  colorlinks=true,
  linkcolor=black,
  citecolor=black,
  urlcolor=blue,
}

\newcommand{\hbarc}{\hbar c}
\newcommand{\Runiv}{R_{\mathrm{universe}}}
\newcommand{\Muniv}{M_{\mathrm{universe}}}
\newcommand{\Nuniv}{N_{\mathrm{universe}}}
\newcommand{\av}[1]{\left\langle #1 \right\rangle}

\title{An Entanglement Mechanism for the Mach--Sciama Relation:\\
Newton's Constant from Cosmic Boundary Conditions}
\author{}
\date{}

\begin{document}

\maketitle

\begin{abstract}
The relation $G \sim c^2 R/M$ between Newton's gravitational constant and
cosmic boundary conditions is well-known in the Mach--Sciama tradition: in
any flat critical-density cosmology, $G c^2 = (\text{const}) \cdot c^4 R/M$
follows from the Friedmann equation. What is \emph{not} explained in standard
cosmology is \emph{why} the universe should sit at this self-consistency point.
We propose a quantum entanglement mechanism --- orthogonal jerk creating
helical trajectories, with bond energy $\hbar c/r$ and torque geometry from
the Bisognano--Wichmann theorem --- that predicts the specific prefactor
$\gamma = 4/5$ from the modular Hamiltonian of the cosmic reduced density
matrix. The relation
\begin{equation}
  G \;=\; \frac{4}{5}\,\frac{c^2 \Runiv}{\Muniv}
\end{equation}
gives a value consistent with the measured $G$ to within current cosmic
parameter uncertainties.

\textbf{Scope.} This is not a free derivation of $G$. The cosmic inputs
$\Runiv$ and $\Muniv$ are themselves determined within general-relativistic
cosmology and contain $G$ implicitly (Section~\ref{sec:scope}). The
substantive content is the prediction of the prefactor $4/5$ from the
modular Hamiltonian, and the mechanism itself, not the numerical value
of $G$. A cleaner $G$-independent test of the same mechanism, the exact
derivation of the Bohr magneton at the electron scale, is given in a
companion paper.
\end{abstract}

\section{Theoretical Foundation}

The derivation rests on five independent results from quantum field theory
and conformal field theory.

\subsection{Bisognano--Wichmann Theorem (1975/1976)}

The modular Hamiltonian for a half-space in relativistic QFT is the generator
of Lorentz boosts perpendicular to the entangling surface:
\begin{equation}
  K = 2\pi \int_{x>0} \mathrm{d}^{d-1}x \; x \; T_{00}(x)
  \label{eq:bw}
\end{equation}
This is a \textbf{theory-independent} result. It applies to any relativistic QFT,
interacting or free. The modular Hamiltonian is always a weighted integral of
the stress-energy tensor $T_{00}$, with the weight being the distance $x$ from
the entangling surface.

\textbf{Physical meaning:} Entanglement creates a boost-rotation geometry, not a
time-translation. The nucleon's spin faces orthogonally to the displacement
toward the partner.

\subsection{Casini--Huerta--Myers (2011)}

For a spherical region of radius $R$ in a $d$-dimensional CFT, the modular
Hamiltonian is:
\begin{equation}
  K = 2\pi \int_{|\vec{x}|<R} \mathrm{d}^{d-1}x \; \beta(r) \; T_{00}(\vec{x})
  \label{eq:chm}
\end{equation}
where the weight function is:
\begin{equation}
  \beta(r) = \frac{R^2 - r^2}{2R}
  \label{eq:chm_weight}
\end{equation}
This is a parabola that vanishes at the boundary $r = R$ and is maximal at
the center $r = 0$. The result is derived by conformally mapping the ball to
a half-space and applying the BW theorem.

\subsection{Cardy--Tonni (2016)}

For an interval of length $l$ in 1+1D CFT, the entanglement Hamiltonian is
derived via conformal mapping from a punctured strip to an annulus:
\begin{equation}
  H_E = 2\pi \int_0^l \mathrm{d}x \; \beta(x) \; T_{00}(x)
  \label{eq:cardy_tonni}
\end{equation}
where the inverse temperature profile is:
\begin{equation}
  \beta(x) = \frac{l}{\pi}
  \frac{\cos(\pi x_0/l) - \cos(\pi x/l)}{\sin(\pi x_0/l)}
  \label{eq:cardy_tonni_beta}
\end{equation}
For the full interval ($x_0 = l/2$), $\cos(\pi/2) = 0$ and $\sin(\pi/2) = 1$, so:
\begin{equation}
  \beta(x) = -\frac{l}{\pi} \cos(\pi x/l)
  \label{eq:cardy_tonni_full}
\end{equation}
The energy scale is set by the inverse of the interval length: $E \sim \hbar c/l$.

\subsection{Bernard et al. (2024)}

The modular Hamiltonian for inhomogeneous free-fermion chains has a commuting
operator $T_A$ whose eigenvalues approximate the entanglement spectrum:
\begin{equation}
  \epsilon_k = \frac{2\pi \lambda_k}{\mathrm{scale} \cdot N^2}
  \label{eq:bernard_eigen}
\end{equation}
For the Racah chain (most general case), the full inverse temperature
$\beta(x) \propto x(L-x)$ is parabolic. When the filling fraction is low
($\rho \ll 1/2$), the Fermi velocity vanishes outside a narrow active region
$[x_-, L]$ near the boundary. Setting $u = (L-x)/L \in [0, 1-x_-/L]$ and
re-expressing $\beta$ on the active interval, the dominant term is
$\beta(r) \propto (1-r/R)^2$ to leading order in the ratio $(L-x_-)/L$
\citep{bernard2024}:
\begin{equation}
  \beta(r) = (1 - r/R)^2
  \label{eq:racah_weight}
\end{equation}

\subsection{Huerta \& van der Velde (2023--2024)}

The modular Hamiltonian for a spherical region in $d$-dimensional CFT can be
obtained by dimensional reduction to the radial semi-infinite line. When a
$d$-dimensional free massless field is decomposed into angular modes, the
modular Hamiltonian decomposes as:
\begin{equation}
  K = \sum_{\ell, \vec{m}} K_{\ell, \vec{m}}
  \label{eq:huerta_decomp}
\end{equation}
where each mode contribution $K_{\ell, \vec{m}}$ is \textbf{exactly} the dimensional
reduction of the parent CFT modular Hamiltonian. The weight function $\beta(r)$
is \textbf{preserved} under dimensional reduction.

Critically, the 1D reduced theories are \textbf{non-conformal} (they carry a
$1/r^2$ potential from the angular Laplacian), yet their modular Hamiltonian
remains local in the energy density with the \textbf{same weight function} as the
$d$-dimensional CFT. This is proven for both scalar fields and Dirac fermions
\citep{huerta2023a, huerta2023b}.

\section{Circularity and Scope}
\label{sec:scope}

Before presenting the mechanism it is essential to be explicit about
what this calculation does and does not establish.

\subsection{The Mach--Sciama Relation Is Not Novel}

In a spatially flat universe at critical density, the Friedmann equation
gives $\rho_c = 3 H^2/(8\pi G)$. The total mass-energy in the Hubble
volume is then
\begin{equation}
  M_H \;=\; \rho_c \cdot \frac{4}{3}\pi R_H^3
       \;=\; \frac{c^2 R_H}{2 G},
       \qquad R_H = c/H,
\end{equation}
which immediately rearranges to $G = c^2 R_H/(2 M_H)$. A relation of the
form $G \sim c^2 R/M$ is therefore present in any standard cosmology;
it has been discussed since Sciama (1953) and is the structural premise
underlying Brans--Dicke theory and several modern emergent-gravity
programs. \emph{Reproducing this structural relation is not, by itself,
new physics.}

\subsection{The Inputs Are Not $G$-Independent}

In standard cosmology, both $\Runiv$ and $\Muniv$ are determined within
the GR/$\Lambda$CDM framework and depend on $G$:

\begin{itemize}
  \item $\Runiv$ is the proper distance to the particle horizon, computed
    from a Friedmann integral that contains $G$ through the matter
    density. Stellar age estimates (white dwarf cooling, globular cluster
    main-sequence turn-off) provide a partially $G$-independent age
    constraint, but the radius $\Runiv = c \cdot t_{\mathrm{age}}$ is at
    best a kinematic proxy.
  \item $\Muniv$ in $\Lambda$CDM is set by the closure condition
    $\Omega_{\mathrm{tot}} = 1$, which contains $G$ explicitly. Dark
    matter mass is determined by gravitational dynamics, so it cannot
    be measured independently of $G$. Only the baryonic component
    ($\Omega_b \approx 0.049$, from BBN and the CMB) is constrained
    by non-gravitational physics.
\end{itemize}

The numerical match between $\frac{4}{5} c^2 \Runiv/\Muniv$ and the
measured $G$ is therefore \emph{not} evidence that the entanglement
mechanism predicts the correct value of $G$ from $G$-free inputs. It
is evidence that the entanglement mechanism is consistent with the
Mach--Sciama relation to within the prefactor.

\subsection{What Is and Is Not Claimed}

\textbf{What is claimed:}
\begin{enumerate}
  \item There is a physically motivated entanglement mechanism whose
    cosmic application produces a relation of the form $G = \gamma\,
    c^2 \Runiv/\Muniv$.
  \item The prefactor $\gamma = 4/5$ is derived from the modular
    Hamiltonian of free fermions on a Racah chain at low filling
    fraction, applied via dimensional reduction to the 3+1D cosmic
    reduced density matrix.
  \item The same mechanism, applied at the electron scale where its
    inputs are $G$-independent, exactly reproduces the Bohr magneton
    (companion paper).
\end{enumerate}

\textbf{What is not claimed:}
\begin{enumerate}
  \item This is not a derivation of $G$ from $G$-free inputs; the
    cosmic parameters used are not independent of $G$.
  \item The $0.21\%$ point-estimate match should not be read as a
    precision result, since $\Runiv$ and $\Muniv$ each carry $5$--$10\%$
    uncertainty in their own right.
  \item The mechanism does not, on its own, replace general relativity.
    What it offers is a candidate explanation for \emph{why} the
    Mach--Sciama relation should hold, with a specific predicted
    prefactor.
\end{enumerate}

The cleanest test of the mechanism is therefore not the gravity
calculation, but the Bohr magneton derivation at the electron scale,
which uses no $G$ at any step.

\section{The Bond Energy $E_{\mathrm{bond}} = \hbar c / r$}

\subsection{Dimensional Analysis}

The modular Hamiltonian $K$ is dimensionless (it appears in $\rho_A = e^{-K}/Z$).
The weight $\beta$ has dimensions $[L]$ (boost parameter = proper distance).
The energy density $T_{00}$ has dimensions $[E]/[L]^d$. The volume element
$\mathrm{d}^d x$ has dimensions $[L]^d$.

\begin{equation}
  K = 2\pi \int \mathrm{d}^d x \; \beta(x) \; T_{00}(x)
  \label{eq:dim_analysis}
\end{equation}

The integral has dimensions $[L]^d \cdot [L] \cdot [E]/[L]^d = [E] \cdot [L]$.
For $K$ to be dimensionless:
\begin{equation}
  \frac{[E] \cdot [L]}{\hbar c} = \mathrm{dimensionless}
  \implies E \sim \frac{\hbar c}{L}
  \label{eq:scaling}
\end{equation}
This is the universal scaling: the entanglement energy between two subsystems
at separation $L$ is $\hbar c/L$, up to a dimensionless coefficient.

\subsection{Why Free Fermions Apply to Nucleons}

The Bernard et al. calculation works with free fermions. Nucleons interact via
the strong force. Why does the free fermion result apply?

\begin{enumerate}
  \item \textbf{BW theorem is theory-independent:} The modular Hamiltonian structure
    $K = 2\pi \int \beta T_{00}$ applies to any relativistic QFT, interacting
    or free. The specific form of $T_{00}$ does not change the local structure.

  \item \textbf{Dimensional analysis is theory-independent:} The scaling
    $E \sim \hbar c/L$ depends only on the dimensions of $T_{00}$ and $\beta$,
    which are universal.

  \item \textbf{Bernard et al. is a verification, not an assumption:} It confirms
    the universal structure in a concrete solvable model.

  \item \textbf{EFT argument:} At cosmic scales, only massless degrees of freedom
    matter. Detailed QCD physics is integrated out and appears only in the
    normalization, not in the $1/r$ scaling.

  \item \textbf{Cross-check:} At the nucleon scale ($r \sim 1$ fm),
    $\hbar c/r \sim 200$ MeV, matching $\Lambda_{\mathrm{QCD}}$. This agreement at
    both nucleon and cosmic scales is independent confirmation.
\end{enumerate}

\subsection{Numerical Verification}

\begin{table}[H]
  \centering
  \caption[Bond energy across scales]{Bond energy $E_{\mathrm{bond}} = \hbar c/r$ across scales.}
  \label{tab:bond_energy}
  \begin{tabular}{llll}
    \toprule
    Scale & $r$ & $E_{\mathrm{bond}} = \hbar c/r$ & Physical interpretation \\
    \midrule
    Nucleon & 1 fm & $3.2 \times 10^{-11}$ J (0.2 MeV) & Strong force scale ($\Lambda_{\mathrm{QCD}}$) \\
    Earth & 6371 km & $5.0 \times 10^{-33}$ J & $3.1 \times 10^{-14}$ eV \\
    Cosmic & $4.4 \times 10^{26}$ m & $7.2 \times 10^{-53}$ J & Gravitational binding scale \\
    \bottomrule
  \end{tabular}
\end{table}

\section{The Geometric Factor $\gamma = 4/5$}

\subsection{The 1D Racah Chain Result}

The cosmic filling fraction is $\rho = \Omega_b = 0.049$ (the baryon fraction).
In the Racah chain framework, this maps to parameters:
\begin{equation}
  x_1 = \frac{\Omega_{\mathrm{DM}}}{\Omega_b} = 5.47, \quad
  x_2 = \frac{\Omega_{\mathrm{DE}}}{\Omega_b} = 13.94, \quad
  x_3 = \frac{\Omega_{\mathrm{DM}}}{\Omega_{\mathrm{DE}}} = 0.39
  \label{eq:racah_params}
\end{equation}
The Fermi velocity $v_F(x)$ is non-zero only on the active region
$[x_-, x_+] = [0.868, 1.000]$, a narrow band near the boundary.
The modular Hamiltonian weight on this region is quadratic:
\begin{equation}
  \beta_{\mathrm{1D}}(x) = (1 - x/L)^2
  \label{eq:racah_1d}
\end{equation}

\subsection{The 3+1D Extension}

The Huerta \& van der Velde theorem proves that the 1D modular Hamiltonian
on the radial line is the dimensional reduction of the $d$-dimensional sphere.
The weight function $\beta(r)$ is \textbf{preserved} under dimensional reduction.

In 3+1D, the spherical volume element $r^2 \mathrm{d}r$ modifies the weighted average.
The weighted average of $1/r$ with the quadratic weight is:
\begin{equation}
  \av{\frac{1}{r}}_{\mathrm{3D}}
  = \frac{\int_0^R \beta(r) \cdot r \, \mathrm{d}r}
         {\int_0^R \beta(r) \cdot r^2 \, \mathrm{d}r}
  \label{eq:weighted_avg}
\end{equation}

For $\beta(r) = (1 - r/R)^2$:

\noindent\textbf{Numerator:}
\begin{equation}
  \int_0^R (1 - r/R)^2 \cdot r \, \mathrm{d}r
  = R^2 \int_0^1 u^2(1-u) \, \mathrm{d}u
  = R^2 \left(\frac{1}{3} - \frac{1}{4}\right)
  = \frac{R^2}{12}
  \label{eq:numerator}
\end{equation}

\noindent\textbf{Denominator:}
\begin{equation}
  \int_0^R (1 - r/R)^2 \cdot r^2 \, \mathrm{d}r
  = R^3 \int_0^1 u^2(1-u)^2 \, \mathrm{d}u
  = R^3 \left(\frac{1}{3} - \frac{1}{2} + \frac{1}{5}\right)
  = \frac{R^3}{30}
  \label{eq:denominator}
\end{equation}

\noindent\textbf{Result:}
\begin{equation}
  \av{\frac{1}{r}} = \frac{R^2/12}{R^3/30} = \frac{5}{2R}
  \label{eq:avg_result}
\end{equation}

The geometric factor is:
\begin{equation}
  \gamma = \frac{2}{\av{1/r} \cdot R} = \frac{2}{5/2} = \frac{4}{5}
  \label{eq:gamma}
\end{equation}

\subsection{Sensitivity to the Weight Exponent}

The quadratic weight $a = 2$ is determined by the Racah chain physics
(low filling fraction $\rho = 0.049$), not by fitting to $G$. A scan
of nearby exponents shows the result is sensitive: the error grows
rapidly as $a$ deviates from 2.

\begin{table}[H]
  \centering
  \caption[Sensitivity to weight exponent]{Sensitivity of $\gamma$ to the weight exponent $a$ in $\beta(r) = (1-r/R)^a$.}
  \label{tab:sensitivity}
  \begin{tabular}{cccc}
    \toprule
    $a$ & $\av{1/r}\,R$ & $\gamma$ & Error vs measured $G$ \\
    \midrule
    1.8 & 2.400 & 0.8333 & $4.39\%$ \\
    1.9 & 2.450 & 0.8163 & $2.26\%$ \\
    \textbf{2.0} & \textbf{2.500} & \textbf{0.8000} & \textbf{$0.21\%$} \\
    2.1 & 2.550 & 0.7843 & $1.75\%$ \\
    2.2 & 2.600 & 0.7692 & $3.64\%$ \\
    \bottomrule
  \end{tabular}
\end{table}

The sensitivity confirms that the quadratic weight is not an arbitrary
choice: a small deviation from $a = 2$ would produce a noticeably worse
agreement with the measured value of $G$.

\section{The Mechanism}

\subsection{Cosmic Entanglement Structure}

In the cosmic reduced density matrix, every nucleon is entangled with every
other nucleon. The angular momentum budget per nucleon is:
\begin{equation}
  s = \frac{\sqrt{3}}{2} \hbar
  \label{eq:spin_budget}
\end{equation}
This budget is distributed across $\Nuniv$ partners, with the
allocation following the modular Hamiltonian structure ($\hbar c/r$).

\subsection{Torque Geometry}

The BW theorem shows that the modular flow is a boost rotation, not a time
translation. The nucleon's spin faces orthogonally to the displacement toward
the partner. This is the ``torque geometry'' of entanglement.

\subsection{Unresolved Entanglement}

When the per-pair coupling between two particles is far below unity, the pair
has not collapsed to a definite correlation. The entanglement amplitude still
exists but remains unresolved. For gravity, this coupling is
$\alpha_{\mathrm{grav}} \approx 5.9 \times 10^{-39}$, so essentially all
nucleon pairs are in this unresolved regime.

This is not zero force. The unresolved amplitude carries the cosmic
entanglement field: each nucleon is entangled with $N_{\mathrm{universe}}$
other nucleons. The balance of angular momentum is the rest of the universe.
In QFT language, this is the reduced density matrix obtained by tracing out
the complement. The universe's isotropic entanglement field is the baseline;
any local mass breaks the isotropy.

\subsection{Gravity as Anisotropy}

The boost rotation back toward the mass is the gravitational force. Gravity
is not a fundamental force---it is the anisotropy in the cosmic entanglement
field created by local mass concentrations.

\section{The Derivation}

\textbf{Terminology note.} Throughout this section we use the symbol $C$ to
denote the per-pair entanglement coupling (the strength with which two
nucleons share entanglement). This is \emph{not} the Wootters concurrence
of two-qubit quantum information theory. The two should not be conflated.

\subsection{Setup}

Each nucleon carries angular momentum $s = \sqrt{3}/2 \cdot \hbar$. From
the modular Hamiltonian (Sections~1--2), the bond energy between two
nucleons at separation $r$ is:
\begin{equation}
  E_{\mathrm{bond}} = \kappa \frac{\hbar c}{r}
  \label{eq:bond_energy}
\end{equation}
where $\kappa$ is an order-unity geometric coefficient. The BW theorem
sets the torque geometry: the nucleon faces orthogonally to the
displacement toward the partner. The isotropic cosmic background
cancels, leaving only the anisotropy from local mass concentrations.

\subsection{Per-Pair Coupling from Energy Balance}

The per-pair coupling $C$ is determined by the total entanglement
energy of the universe. The universe's mass-energy is stored as
entanglement energy across all nucleon pairs:
\begin{equation}
  E_{\mathrm{entanglement}}
  = \frac{\Nuniv^2}{2} \cdot C \cdot \kappa \hbar c \cdot \av{\frac{1}{r}}
  \label{eq:total_entangle}
\end{equation}
where $\av{1/r} = 5/(2\Runiv)$ from the quadratic weight.
Equating to the universe's mass-energy $\Muniv c^2$:
\begin{equation}
  \frac{N^2}{2} \cdot C \cdot \kappa \hbar c \cdot \frac{5}{2R} = M c^2
  \label{eq:energy_balance}
\end{equation}
Solving for $C$ and substituting $N = M/m_p$:
\begin{equation}
  C = \frac{4}{5} \frac{m_p^2 c \Runiv}{\kappa \Muniv \hbar}
  \label{eq:concurrence}
\end{equation}

\subsection{$G$ Emerges}

The force per nucleon pair is:
\begin{equation}
  F_{\mathrm{pair}} = C \cdot \kappa \frac{\hbar c}{r^2}
  \label{eq:pair_force}
\end{equation}
The total force between two masses $M$ and $m$ at distance $r$:
\begin{equation}
  F = \frac{M}{m_p} \cdot \frac{m}{m_p} \cdot C \cdot \kappa \frac{\hbar c}{r^2}
  \label{eq:total_force}
\end{equation}
Substituting \cref{eq:concurrence}:
\begin{equation}
  F = \frac{4}{5} \cdot M \cdot m \cdot \frac{c^2 \Runiv}{\Muniv} \cdot \frac{1}{r^2}
  \label{eq:newton_form}
\end{equation}
Comparing to Newton's law $F = G M m / r^2$:
\begin{equation}
  G = \frac{4}{5} \frac{c^2 \Runiv}{\Muniv}
  \label{eq:final_g}
\end{equation}
\textbf{Note:} $\hbar$, $m_p$, and $\kappa$ all cancel. The final expression depends
only on $c$, $\Runiv$, and $\Muniv$.

\section{Numerical Consistency Check}

Using cosmic inputs:
\begin{align}
  c &= 2.998 \times 10^8 \; \mathrm{m/s} \\
  \Runiv &= 4.4 \times 10^{26} \; \mathrm{m} \\
  \Muniv &= 4.73 \times 10^{53} \; \mathrm{kg}
\end{align}
the predicted relation gives
\begin{equation}
  G_{\mathrm{predicted}}
  = \frac{4}{5} \frac{(2.998 \times 10^8)^2 \cdot (4.4 \times 10^{26})}
                     {4.73 \times 10^{53}}
  = 6.688 \times 10^{-11} \; \mathrm{m^3/(kg{\cdot}s^2)},
  \label{eq:g_derived}
\end{equation}
to be compared with $G_{\mathrm{measured}} = 6.674 \times 10^{-11}\;
\mathrm{m^3/(kg{\cdot}s^2)}$. The ratio is $1.0021$.

This check should be read in the light of Section~\ref{sec:scope}: the
inputs $\Runiv$ and $\Muniv$ are not independent of $G$, and each carries
roughly $5$--$10\%$ uncertainty in standard cosmology. The headline
$0.21\%$ agreement is therefore a consistency check, not a precision
result. The honest statement is: the predicted prefactor $4/5$, combined
with the standard cosmic parameters, is consistent with the measured value
of $G$ to well within those parameter uncertainties.

A genuine prefactor test would require an independent constraint on
$\Muniv/\Runiv$ that does not pass through $G$. We are not aware of one
that is currently tight enough to discriminate, say, $4/5$ from $1/2$ or
$1$ at the few-percent level.

\section{Status of the Predictions}

Honest accounting of what the mechanism currently predicts cleanly,
predicts structurally, and does not yet predict.

\subsection{Cleanly Predicted (G-Free): Bohr Magneton}

Applied at the electron scale, the same helical-jerk mechanism predicts
\begin{equation}
  \mu_B = \frac{e \hbar}{2 m_e}
  \label{eq:bohr_magneton_ref}
\end{equation}
exactly, with no fitted parameters and using only $e$, $\hbar$, $m_e$,
and $c$ --- none of which involve $G$. This is the cleanest test of the
mechanism and is treated in detail in the companion paper.

\subsection{Structurally Predicted (Quantitatively Off): Frame Dragging}

The boost generator structure (BW theorem) implies that rotating masses
drag the local entanglement field, producing a Lense--Thirring-like
effect. The mechanism predicts the correct \emph{structure}: frame
dragging is a boost effect of the modular flow under rotation, not a
separate phenomenon.

The \emph{quantitative} match is currently off by roughly a factor of
$10^3$. A naive collective treatment (random walk of $\sqrt{N_{\mathrm{Earth}}}$
nucleons) does not reproduce $\omega_{\mathrm{FD}}^{\mathrm{GR}} =
2GJ/(c^2 R^3) \approx 1.9 \times 10^{-14}$ rad/s. Closing this gap
requires a more precise treatment of the collective contribution at
the nucleon level, and is open work.

\subsection{Order-of-Magnitude: Dark Energy}

The unresolved entanglement field at the cosmic scale gives an energy
density of order $\rho_{\mathrm{DE}} \approx 8.1 \times 10^{-11}$ J/m$^3$,
compared to the observed $\rho_{\mathrm{DE}}^{\mathrm{obs}} \approx
6.9 \times 10^{-10}$ J/m$^3$. The prediction is within an order of
magnitude. The factor of $\sim 8$ is consistent with the holographic
picture (dark energy stored on the cosmic horizon, a 2D surface,
rather than distributed through the bulk volume), but a precise
holographic accounting is beyond the present scope.

\subsection{Not a Free Prediction: $\alpha_{\mathrm{grav}}$}

The standard gravitational coupling
$\alpha_{\mathrm{grav}} = G m_p^2/(\hbar c) \approx 5.9 \times 10^{-39}$
follows from the predicted relation for $G$ together with $m_p$, $\hbar$,
and $c$. It is therefore consistent with the framework but not an
independent prediction --- it inherits the same circularity as $G$ itself.

\section{Comparison to Kolmogorov's Four-Fifths Law}

The factor $4/5$ also appears as an exact coefficient in Kolmogorov's
four-fifths law for fully developed turbulence, where the third-order
longitudinal velocity structure function satisfies
$\langle (\delta u_\parallel)^3 \rangle = -\frac{4}{5} \varepsilon r$
in the inertial range. This is one of the few exact results in
turbulence theory, deriving from energy conservation in a cascade
across scales. The recurrence of $4/5$ in two distinct
energy-transfer-across-scales problems (turbulent cascade in 3D,
entanglement-energy distribution across the cosmic 3+1D modular
Hamiltonian) is suggestive but not by itself a derivation; we note
the parallel as a possible structural connection.

\section{Computational Verification}

All derivations are implemented in Python scripts:

\begin{itemize}
  \item \texttt{derive\_G.py}---Main $G$ derivation (752 lines)
  \item \texttt{derive\_bond\_energy.py}---$E_{\mathrm{bond}} = \hbar c/r$ derivation (420 lines)
  \item \texttt{derive\_3d\_extension.py}---3+1D extension via Huerta \& van der Velde (350 lines)
  \item \texttt{constants.py}---Physical constants and helper functions
\end{itemize}

All scripts run cleanly and produce consistent results.

\bibliographystyle{plain}
\bibliography{references}

\end{document}
